\DocumentMetadata{
    tagging=on,
    tagging-setup={math/setup=mathml-SE},
    pdfstandard=ua-2,
    lang=en-US
}
\documentclass[twoside,12pt]{article}

\usepackage{tikzStickFigure}
\usepackage{unicode-math}

\usepackage[outer=0.75in,inner=1in]{geometry}
\usepackage{enumitem}
\setlist[description]{
    % style=nextline,
    labelwidth=0pt,
    leftmargin=0.5in,
    itemindent=-12pt,
}
\usepackage{geometry}
\geometry{inner=0.75in,outer=0.5in,top=1in,bottom=0.75in}


\title{Tikz Stick Figures}
\author{C.F. Rocca Jr.}



\begin{document}

\maketitle

\begin{abstract}
    Tikz Stick Figures provides options for drawing basic stick figures to add to documents.  The figures may be equipped with a variety of accessories including headgear, weapons, and clothing.  Most of the options were inspired by figures that could accompany documents related to Tabletop Role Playing Games (TTRPGs).
\end{abstract}

\section{Sampler}

\begin{tikzpicture}[scale=0.5,transform shape]
%% Basic stick figure with nose at (0,0)
    \stickfig[];

%%% Basic Elf
    \stickfig[
        position = {(5,0)},
        species = elf,
        hands = down,
        weapon = spear,
    ];

%%% Dwarf Barbarian
    \stickfig[
        position = {(10,0)},
        species = dwarf,
        hands = fight left,
        stance = run left,
        weapon = axe,
        weapon angle = 100,
        weapon hand = left hand,
        hat = helmet,
        hat color = red,
    ];

%%% Human Fighter
    \stickfig[
        position = {(15,0)},
        species = human,
        hands = fight right,
        stance = run right,
        weapon = sword,
        weapon angle = 80,
        shield = True,
        shield color = green!50!black,
        star color = pink,
    ];

%%% Casualty
\begin{scope}[rotate around = {70:(18,-1)}]
    \stickfig[
        position = {(18,-1)},
        condition = injured,
    ]
\end{scope}

%%% Fully Decked Out Named Wizard
    \stickfig[
        stance=run left,position={(20,-10)},name=Chuck,hat=wizard hat,
        weapon=bow,weapon hand=left hand,weapon angle =100,
        shield=True, shield hand=right hand,hat color=blue,
        condition = injured,species=elf
    ];
    \node[font=\large\bfseries,anchor=north] at (\name right foot) {\name};

    % \stickfig[species = dwarf, condition=injured]
    % \stickfig[species = elf, condition=injured, weapon = axe, weapon angle = 0]
    % \stickfig[species = human, condition=injured]

    % \path[line width = 2,fill=brown,bend angle =80] ($(face)+(1.05,0)$) 
    %     arc (0:-73:1.5)
    %     arc (253:180:1.5)
    %     to[bend right] ($(face)+(0,-10pt)$)
    %     to[bend right] ($(face)+(1.05,0)$);


%% Unnamed stick figure 
    % \stickfig[
    %     stance={run right},position={(1,5)},shield=True,weapon=axe,
    %     condition=injured,weapon angle=-35,species=elf,hands = going right,
    % ];

    % \stickfig[scale=0.75,position={(-4,6)},hands=up]
    % \stickfig[scale=0.75,position={(0,6)},hands=down]

%% Combatents
    % \stickfig[
    %     scale=0.66,position={(4,6)},hands=fight left,
    %     shield=True,shield hand=right hand,,shield color=blue,
    %     weapon=axe,weapon hand=left hand,weapon angle=0,
    %     hat=helmet,hat color=pink,
    %     stance=run left,
    %     condition = injured
    % ]
    % \stickfig[
    %     scale=0.75,position={(8,6)},hands=fight right,
    %     shield=True,
    %     weapon=sword,weapon angle=70,
    %     hat=wizard hat,hat color=purple,
    %     stance=run right,
    % ]

%% Named stick figure
    % \stickfig[
    %     stance=run left,position={(5,-3)},name=chuck,hat=wizard hat,scale=0.75,
    %     weapon=bow,weapon hand=left hand,weapon angle =100,
    %     shield=True, shield hand=right hand,hat color=blue,
    % ];

%% Sample Lone Accessories
    % \pic at (-3,-3) {shield={-3,-3}{green}{yellow}};
    % \pic at (-3,-1) {helmet={-3,-3}{black}{yellow}};
    % \pic at (-3,1) {axe={-3,0}{0}};

%% Helper Layout Grid
    \draw[thin,dashed] (0,0) grid (20,-10);

%% Development Space

\end{tikzpicture}

\clearpage

\section{Stick Figures}
\subsection{Stick Figure Basic Examples}

\subsection{Stick Figure Key Values}

A stick figure can be added with just the \verb|\stickfigure| command which produces a standing figure with raised arms.  But there are a lot of options for making out stick figures more interesting. Figure \ref{fig:figure_options} gives a comprehensive sample of all the possible options; lists of options begin with the default. 

\begin{figure}[hpbt]
\begin{verbatim}
\stickfig[
    scale = <percent scale 0 ...>,
    position = {(x,y)},
    name = <figures name>,
    condition = <healthy, injured>,
    species = <human, elf, dwarf>,
    stance = <standing, run right, run left>,
    hands = <up, down, bent right, bent left, fight right, fight left>,
    hat = <none, wizard hat, helmet>,
    hat color = <color>,
    weapon = <none, sword, spear, bow, axe>, 
    weapon hand = <right hand, left hand>,
    weapon angle = <rotation angle>,
    shield = <False, True>,
    shield hand = <left hand, right hand>,
    shiekd color = <color>,
    star color = <color>,
];
\end{verbatim}
    \caption{Stick Figure Options and Syntax}
    \label{fig:figure_options}
\end{figure}


\begin{description}
    \item[condition] injured or not
    \item[hands] \verb|<up>|, \verb|<down>|, \verb|<bent right>|, \verb|<bent left>|, \verb|<fight right>|, \verb|<fight left>|
    \item[hat] \verb|<none>|, \verb|<wizard hat>|, \verb|<helmet>|
    \item[hat color] ...
    \item[name] ...
    \item[position] ...
    \item[scale] ...
    \item[shied] shield true or false
    \item[shield color] ...
    \item[shield hand] ...
    \item[species] ...
    \item[stance] \verb|<standing>|, \verb|<run left>|, \verb|<run right>|
    \item[weapon] ...
    \item[weapon angle] ...
    \item[weapon hand] ...
\end{description}

\section{Accessories}

\begin{description}
    \item[Axe] \tikz{\pic[scale=0.5] at (0,0){axe};}
    \item[Bow] \tikz{\pic[scale=0.5] at (0,0){bow};}
    \item[Helmet] \tikz{\pic[scale=0.5] at (0,0){helmet};}
    \item[Shield] \tikz{\pic[scale=0.5] at (0,0){shield};}
    \item[Spear] \tikz{\pic[scale=0.5] at (0,0){spear};}
    \item[Sword] \tikz{\pic[scale=0.5] at (0,0){sword};}
    \item[Wizard's Hat] \tikz{\pic[scale=0.5] at (0,0){wizardhat};}
\end{description}



\end{document}
